% ============================================================
%  第一卷 第二章 守恒定律（§6–§10 + 习题）
%  OCR 错误修正: ι→i, \hat{\otimes}→δ, \mathfrak{z}→z, ρ→p,
%  习题3七处中文括注按 Landau 原书内容重建; §10 力学相似性于全书装配时补排,
%  习题1答案 OCR 漏撇号 √(m/m)→√(m′/m) 按原书修正
% ============================================================
\chapter{守恒定律}

\section{能量}

在力学系统运动过程中，描述其状态的 $2s$ 个变量 $q_i,\dot{q}_i\ (i=1,2,\cdots,s)$ 随时间变化.但是存在关于这些变量的某些函数，其值在运动过程中保持恒定，且仅由初始条件决定，这样的函数称为运动积分.

对于具有 $s$ 个自由度的封闭力学系统，独立的运动积分数等于 $2s-1$，这从下面简单的论证中很容易得出.运动方程通解中包含 $2s$ 个任意常数（参见方程 (2.6) 下面的讨论）.由于封闭系统的运动方程不显含时间，所以可以完全任意选择初始时刻，总可以将方程解中的任意常数之一选作时间的可加常数 $t_0$.从 $2s$ 个函数
\begin{equation*}
\begin{aligned}
  q_i&=q_i\left(t+t_0,C_1,C_2,\cdots,C_{2s-1}\right),\\
  \dot{q}_i&=\dot{q}_i\left(t+t_0,C_1,C_2,\cdots,C_{2s-1}\right),
\end{aligned}
\end{equation*}
中消去 $t+t_0$，将 $2s-1$ 个任意常数 $C_1,C_2,\cdots,C_{2s-1}$ 表示成 $q_i,\dot{q}_i$ 的函数，这些函数就是运动积分.

然而，并不是所有的运动积分在力学中有相同的重要性，其中一些运动积分源自时间和空间的均匀性和各向同性这样的基本性质，这种运动积分的恒定不变性才具有深刻的意义.由这些运动积分表示的量是守恒量，它们具有一个重要的共同性质——可加性.对于几个相互独立部分组成的系统，守恒量的值等于各个部分相应值之和.

可加性使相应的物理量在力学中具有特别重要的作用.例如，假设两个物体在某时间间隔内相互作用.由于在相互作用发生前后，整个系统的每个可加运动积分值就等于两物体相应值之和，如果已知作用发生前各个物体的运动状态，那么通过守恒定律就能立即得到与相互作用发生后各物体运动状态有关的各种结论.

我们首先介绍时间均匀性导出的守恒定律.

由于时间具有均匀性，封闭系统的拉格朗日函数不显含时间.因此拉格朗日函数对时间的全导数可以写成
\begin{equation*}
  \frac{\mathrm{d}L}{\mathrm{d}t}
  =\sum_i\frac{\partial L}{\partial q_i}\dot{q}_i
  +\sum_i\frac{\partial L}{\partial\dot{q}_i}\ddot{q}_i
\end{equation*}
（如果 $L$ 显含时间，右端还应该有 $\partial L/\partial t$）.利用拉格朗日方程将 $\partial L/\partial q_i$ 替换为 $\dfrac{\mathrm{d}}{\mathrm{d}t}\dfrac{\partial L}{\partial\dot{q}_i}$，得
\begin{equation*}
  \frac{\mathrm{d}L}{\mathrm{d}t}
  =\sum_i\dot{q}_i\frac{\mathrm{d}}{\mathrm{d}t}\frac{\partial L}{\partial\dot{q}_i}
  +\sum_i\frac{\partial L}{\partial\dot{q}_i}\ddot{q}_i
  =\sum_i\frac{\mathrm{d}}{\mathrm{d}t}\left(\frac{\partial L}{\partial\dot{q}_i}\dot{q}_i\right),
\end{equation*}
或者
\begin{equation*}
  \frac{\mathrm{d}}{\mathrm{d}t}\left(\sum_i\dot{q}_i
  \frac{\partial L}{\partial\dot{q}_i}-L\right)=0.
\end{equation*}
由此可知
\begin{equation}
  E=\sum_i\dot{q}_i\frac{\partial L}{\partial\dot{q}_i}-L
\end{equation}
在封闭系统运动中保持不变，是运动积分，称为系统的能量.根据 (6.1)，能量与拉格朗日函数的关系是线性的，由拉格朗日函数的可加性可以直接得出能量的可加性.

在上述推导中仅仅利用了拉格朗日函数不显含时间的性质，所以能量守恒定律不仅对于封闭系统成立，对位于定常外场（即不显含时间）中的系统也成立.能量守恒的力学系统也称为保守系统.

在§5我们已经知道，封闭（或者位于定常外场中的）系统的拉格朗日函数可写成
\begin{equation*}
  L=T(q,\dot{q})-U(q),
\end{equation*}
其中 $T$ 是速度的二次函数.利用著名的齐次函数的欧拉定理可得
\begin{equation*}
  \sum_i\dot{q}_i\frac{\partial L}{\partial\dot{q}_i}
  =\sum_i\dot{q}_i\frac{\partial T}{\partial\dot{q}_i}=2T.
\end{equation*}
将此式代入 (6.1) 得
\begin{equation}
  E=T(q,\dot{q})+U(q),
\end{equation}
用笛卡儿坐标写成
\begin{equation}
  E=\sum_a\frac{m_a v_a^2}{2}+U(\boldsymbol{r}_1,\boldsymbol{r}_2,\cdots).
\end{equation}
可见，系统的能量可以表示为本质不同的两项之和：依赖于速度的动能和仅仅依赖于质点坐标的势能.

\section{动量}

另一个守恒定律与空间的均匀性相关.

根据空间均匀性，封闭力学系统在空间中整体平移时，其性质保持不变.因此我们研究一个无穷小平移 $\boldsymbol{\varepsilon}$，并求拉格朗日函数保持不变的条件.

平移就是将系统中所有质点移动相同的位移 $\boldsymbol{\varepsilon}$ 的变换，即 $\boldsymbol{r}_a\to\boldsymbol{r}_a+\boldsymbol{\varepsilon}$.在速度不变时，坐标的无穷小的改变使拉格朗日函数产生的变化为
\begin{equation*}
  \delta L=\sum_a\frac{\partial L}{\partial\boldsymbol{r}_a}\cdot\delta\boldsymbol{r}_a
  =\boldsymbol{\varepsilon}\cdot\sum_a\frac{\partial L}{\partial\boldsymbol{r}_a},
\end{equation*}
其中求和是对系统中的所有质点进行的.对任意 $\boldsymbol{\varepsilon}$ 要求 $\delta L=0$ 等价于
\begin{equation}
  \sum_a\frac{\partial L}{\partial\boldsymbol{r}_a}=0.
\end{equation}
根据拉格朗日方程 (5.2) 得
\begin{equation*}
  \sum_a\frac{\mathrm{d}}{\mathrm{d}t}\frac{\partial L}{\partial\boldsymbol{v}_a}
  =\frac{\mathrm{d}}{\mathrm{d}t}\sum_a\frac{\partial L}{\partial\boldsymbol{v}_a}=0.
\end{equation*}
于是我们可得结论：封闭力学系统的矢量
\begin{equation}
  \boldsymbol{P}=\sum_a\frac{\partial L}{\partial\boldsymbol{v}_a}
\end{equation}
在运动中保持不变.矢量 $\boldsymbol{P}$ 称为系统的动量.对拉格朗日函数 (5.1) 求导可得用质点速度表示的动量
\begin{equation}
  \boldsymbol{P}=\sum_a m_a\boldsymbol{v}_a.
\end{equation}
动量的可加性是显然的.与能量不同之处在于，无论质点之间的相互作用是否可以忽略，系统的动量都等于各个质点的动量
\begin{equation*}
  \boldsymbol{p}_a=m_a\boldsymbol{v}_a
\end{equation*}
之和.

只有在没有外场的情况下，动量矢量的三个分量都守恒.然而，在有外场的情况下，如果势能不显含某个笛卡儿坐标，则相应的该方向的动量分量守恒.显然，沿着这个不出现在势能中的坐标相应的坐标轴平移不会改变力学系统的性质，动量在该轴上投影守恒.例如，在方向沿着 $z$ 轴的均匀场中，沿着 $x$ 和 $y$ 轴的动量分量守恒.

等式 (7.1) 的物理含义非常简单.导数 $\partial L/\partial\boldsymbol{r}_a=-\partial U/\partial\boldsymbol{r}_a$ 是作用在第 $a$ 个质点上的力 $\boldsymbol{F}_a$.等式 (7.1) 表明，作用在封闭系统的所有质点上的力之和等于零
\begin{equation}
  \sum_a\boldsymbol{F}_a=0.
\end{equation}
特别地，当系统只由两个质点组成时，$\boldsymbol{F}_1+\boldsymbol{F}_2=0$，两个质点的相互作用力大小相等、方向相反.这就是著名的作用与反作用互等定律（牛顿第三定律）.

如果用广义坐标 $q_i$ 描述运动，则拉格朗日函数对广义速度的导数
\begin{equation}
  p_i=\frac{\partial L}{\partial\dot{q}_i}
\end{equation}
称为广义动量，而它对广义坐标的导数
\begin{equation}
  F_i=\frac{\partial L}{\partial q_i}
\end{equation}
称为广义力.采用上述符号，拉格朗日方程可以写成
\begin{equation}
  p_i=F_i.
\end{equation}
在笛卡儿坐标下广义动量就是矢量 $\boldsymbol{p}_a$ 的分量.一般情况下 $p_i$ 是广义速度 $\dot{q}_i$ 的线性齐次函数，不能化为质量与速度的积.

\begin{xiti}

\noindent{\heiti 习题}\quad 质量为 $m$ 的质点以速度 $\boldsymbol{v}_1$ 从一个势能为常数 $U_1$ 的半空间运动到另一个势能为常数 $U_2$ 的半空间.求质点运动方向的改变.

解：势能不依赖于其轴平行两个半空间分界面的坐标，因此质点动量在该分界面上的投影守恒.用 $\theta_1,\theta_2$ 表示质点穿越分界面前后速度 $\boldsymbol{v}_1,\boldsymbol{v}_2$ 与分界面法线的夹角，于是有 $v_1\sin\theta_1=v_2\sin\theta_2$.$v_1$ 和 $v_2$ 之间的关系由能量守恒定律给出，最后可求得
\begin{equation*}
  \frac{\sin\theta_1}{\sin\theta_2}
  =\sqrt{1+\frac{2}{mv_1^2}(U_1-U_2)}.
\end{equation*}

\end{xiti}

\section{质心}

封闭系统的动量对于不同的惯性参考系有不同的值.如果参考系 $K'$ 相对参考系 $K$ 以速度 $\boldsymbol{V}$ 运动，则质点相对这两个参考系的速度 $\boldsymbol{v}'_a$ 和 $\boldsymbol{v}_a$ 满足关系式 $\boldsymbol{v}_a=\boldsymbol{v}'_a+\boldsymbol{V}$.因此在这两个参考系中动量值 $\boldsymbol{P}$ 和 $\boldsymbol{P}'$ 满足关系式
\begin{equation*}
  \boldsymbol{P}=\sum_a m_a\boldsymbol{v}_a
  =\sum_a m_a\boldsymbol{v}'_a+\boldsymbol{V}\sum_a m_a,
\end{equation*}
或者
\begin{equation}
  \boldsymbol{P}=\boldsymbol{P}'+\boldsymbol{V}\sum_a m_a.
\end{equation}
特别地，一定存在使得总动量等于零的参考系 $K'$.令 (8.1) 中 $\boldsymbol{P}'=0$，求得参考系 $K'$ 的速度为
\begin{equation}
  \boldsymbol{V}=\frac{\boldsymbol{P}}{\sum m_a}
  =\frac{\sum m_a\boldsymbol{v}_a}{\sum m_a}.
\end{equation}
如果在给定参考系下力学系统的总动量为零，则称系统相对该参考系静止.这是单个质点静止概念的自然推广.公式 (8.2) 给出的速度 $\boldsymbol{V}$，具有动量不为零的力学系统“整体运动”速度的含义.由此可见，动量守恒定律自然地给出了系统整体静止和速度的概念.

公式 (8.2) 还表明，动量 $\boldsymbol{P}$ 和系统整体运动速度 $\boldsymbol{V}$ 的关系，就如同一个质点动量和速度的关系，该质点的质量等于系统中所有质点的质量之和 $\mu=\sum m_a$.这正表示了质量的可加性.

公式 (8.2) 右端可以看作是下面表达式对时间的导数
\begin{equation}
  \boldsymbol{R}=\frac{\sum m_a\boldsymbol{r}_a}{\sum m_a}.
\end{equation}
可以说，系统整体运动速度就是径矢为 (8.3) 的点在空间中的运动速度，这个点称为系统的质心.

封闭系统动量守恒定律可以表述为：系统的质心作匀速直线运动.这是§3给出的单一自由质点惯性定律的推广，单一自由质点的质心就是质点本身.

在研究封闭系统的力学性质时，自然采用质心静止的参考系，这就可以不必研究系统整体的匀速直线运动，这样的运动并不重要.

整体静止的力学系统的能量通常称为内能 $E_\mathrm{int}$，它包括系统内质点的相对运动动能和相互作用势能.以速度 $\boldsymbol{V}$ 作整体运动的系统的能量可以写成
\begin{equation}
  E=\frac{\mu V^2}{2}+E_\mathrm{int}.
\end{equation}
尽管这个公式非常显然，但我们还是给出以下推导.力学系统相对参考系 $K$ 和 $K'$ 的能量 $E$ 和 $E'$ 的关系为
\begin{equation*}
\begin{aligned}
  E&=\frac{1}{2}\sum_a m_a v_a^2+U
  =\frac{1}{2}\sum_a m_a\left(\boldsymbol{v}'_a+\boldsymbol{V}\right)^2+U\\
  &=\frac{\mu V^2}{2}+\boldsymbol{V}\cdot\sum_a m_a\boldsymbol{v}'_a
  +\frac{1}{2}\sum_a m_a v_a'^2+U
\end{aligned}
\end{equation*}
或者
\begin{equation}
  E=E'+\boldsymbol{V}\cdot\boldsymbol{P}'+\frac{\mu V^2}{2}.
\end{equation}
这个公式给出了相对两个不同惯性参考系的能量关系，类似于公式 (8.1) 给出的动量关系.如果在参考系 $K'$ 中系统质心静止，则 $\boldsymbol{P}'=0$，$E'=E_\mathrm{int}$，这时就得到公式 (8.4).

\begin{xiti}

\noindent{\heiti 习题}\quad 求相对两个不同惯性参考系的作用量之间的变换关系.

解：拉格朗日函数等于动能和势能之差，显然，它按照类似于公式 (8.5) 的形式变换
\begin{equation*}
  L=L'+\boldsymbol{V}\cdot\boldsymbol{P}'+\frac{\mu V^2}{2}.
\end{equation*}
将该等式对时间积分可得所要求的作用量的变换关系
\begin{equation*}
  S=S'+\mu\boldsymbol{V}\cdot\boldsymbol{R}'
  +\frac{\mu V^2}{2}t,
\end{equation*}
其中 $\boldsymbol{R}'$ 是在参考系 $K'$ 中系统的质心的径矢.

\end{xiti}

\section{角动量}

下面研究由空间各向同性得到的守恒定律.

各向同性意味着封闭系统整体在空间中任意转动时，力学性质保持不变.因此，我们研究系统整体的无穷小转动并求出拉格朗日函数保持不变的条件.

我们引入无穷小转动矢量 $\delta\boldsymbol{\varphi}$，其大小等于转角 $\delta\varphi$，方向沿着转动轴（转动方向与 $\delta\boldsymbol{\varphi}$ 的方向之间符合右手螺旋法则）.

我们首先研究，在系统转动时，从坐标原点（位于转动轴上）指向系统中任意质点的径矢的位移.径矢端点的线位移与转角的关系为（如图5所示）
\begin{equation*}
  |\delta\boldsymbol{r}|=r\sin\theta\cdot\delta\varphi.
\end{equation*}
位移矢量的方向垂直于过 $\boldsymbol{r}$ 和 $\delta\boldsymbol{\varphi}$ 的平面.显然有
\begin{equation}
  \delta\boldsymbol{r}=\delta\boldsymbol{\varphi}\times\boldsymbol{r}.
\end{equation}
在系统转动时不仅径矢的方向改变，而且所有质点的速度也改变，并且所有矢量的变化规律相同.所以速度相对固定坐标系的增量为
\begin{equation}
  \delta\boldsymbol{v}=\delta\boldsymbol{\varphi}\times\boldsymbol{v}.
\end{equation}

\par\nopagebreak\noindent\begin{minipage}{\linewidth}\centering\includegraphics[width=4.5cm]{fig5.pdf}\\[0.3em]{\small 图 5}\end{minipage}\par\nopagebreak

将这些表达式代入转动时拉格朗日函数不变的条件
\begin{equation*}
  \delta L=\sum_a\left(\frac{\partial L}{\partial\boldsymbol{r}_a}
  \cdot\delta\boldsymbol{r}_a
  +\frac{\partial L}{\partial\boldsymbol{v}_a}\cdot\delta\boldsymbol{v}_a\right)=0
\end{equation*}
并做代换 $\partial L/\partial\boldsymbol{v}_a=\boldsymbol{p}_a$，$\partial L/\partial\boldsymbol{r}_a=\dot{\boldsymbol{p}}_a$，得
\begin{equation*}
  \sum_a\left[\dot{\boldsymbol{p}}_a\cdot(\delta\boldsymbol{\varphi}
  \times\boldsymbol{r}_a)
  +\boldsymbol{p}_a\cdot(\delta\boldsymbol{\varphi}
  \times\boldsymbol{v}_a)\right]=0.
\end{equation*}
置换因子的次序并将 $\delta\boldsymbol{\varphi}$ 移到求和号之外
\begin{equation*}
  \delta\boldsymbol{\varphi}\cdot\sum_a
  \left(\boldsymbol{r}_a\times\dot{\boldsymbol{p}}_a
  +\boldsymbol{v}_a\times\boldsymbol{p}_a\right)
  =\delta\boldsymbol{\varphi}\cdot\frac{\mathrm{d}}{\mathrm{d}t}
  \sum_a\boldsymbol{r}_a\times\boldsymbol{p}_a=0.
\end{equation*}
由 $\delta\boldsymbol{\varphi}$ 的任意性可得
\begin{equation*}
  \frac{\mathrm{d}}{\mathrm{d}t}\sum_a\boldsymbol{r}_a\times\boldsymbol{p}_a=0,
\end{equation*}
即在封闭力学系统运动过程中矢量
\begin{equation}
  \boldsymbol{M}=\sum_a\boldsymbol{r}_a\times\boldsymbol{p}_a
\end{equation}
保持不变，这个物理量称为系统的角动量\footnote{也称动量矩.}.类似于线动量，这个物理量不依赖于质点之间是否有相互作用，它的可加性是显然的.

可加的运动积分就这些.就是说，任何封闭系统总共有 7 个这样的运动积分：能量、动量的三个分量和角动量的三个分量.

既然在角动量的定义中出现了质点的径矢，它的取值通常就与坐标原点的选取有关.假定两个坐标原点相差矢量 $\boldsymbol{a}$，同一个点对这两个坐标原点的径矢分别为 $\boldsymbol{r}_a$ 和 $\boldsymbol{r}'_a$，则有关系式 $\boldsymbol{r}_a=\boldsymbol{r}'_a+\boldsymbol{a}$.因此有
\begin{equation*}
  \boldsymbol{M}=\sum_a\boldsymbol{r}_a\times\boldsymbol{p}_a
  =\sum_a\boldsymbol{r}'_a\times\boldsymbol{p}_a
  +\boldsymbol{a}\times\sum_a\boldsymbol{p}_a,
\end{equation*}
即
\begin{equation}
  \boldsymbol{M}=\boldsymbol{M}'+\boldsymbol{a}\times\boldsymbol{P}.
\end{equation}
由此可知，只有在系统整体静止（即 $\boldsymbol{P}=0$）时，其角动量不依赖于坐标原点的选择.角动量值的这种不确定性不会影响到角动量守恒定律，因为封闭系统的动量也守恒.

我们来推导相对不同惯性参考系 $K$ 和 $K'$ 的角动量之间的关系.设参考系 $K'$ 相对 $K$ 的速度为 $\boldsymbol{V}$，假定它们的坐标原点在某给定时刻重合.那么质点相对两个参考系的径矢相同，速度满足关系式 $\boldsymbol{v}_a=\boldsymbol{v}'_a+\boldsymbol{V}$.于是有
\begin{equation*}
  \boldsymbol{M}=\sum_a m_a\boldsymbol{r}_a\times\boldsymbol{v}_a
  =\sum_a m_a\boldsymbol{r}_a\times\boldsymbol{v}'_a
  +\sum_a m_a\boldsymbol{r}_a\times\boldsymbol{V}.
\end{equation*}
右端第一项是相对参考系 $K'$ 的角动量 $\boldsymbol{M}'$，在第二项中利用质心径矢公式 (8.3)，可得
\begin{equation}
  \boldsymbol{M}=\boldsymbol{M}'+\mu\boldsymbol{R}\times\boldsymbol{V}.
\end{equation}
这个公式给出了相对不同参考系的角动量之间的变换关系，与能量关系式 (8.5) 和动量关系式 (8.1) 类似.

如果系统整体相对参考系 $K'$ 静止，则 $\boldsymbol{V}$ 是系统质心的速度，而 $\mu\boldsymbol{V}$ 是系统相对于参考系 $K$ 的总动量 $\boldsymbol{P}$，进而有
\begin{equation}
  \boldsymbol{M}=\boldsymbol{M}'+\boldsymbol{R}\times\boldsymbol{P}.
\end{equation}
就是说，力学系统的角动量由其相对静止的参考系中的“内禀角动量”和整体运动的角动量 $\boldsymbol{R}\times\boldsymbol{P}$ 构成.

虽然只有封闭系统的角动量（对任意坐标原点）三个分量都守恒，但是在一定限制下，这个守恒定律对于在外场中运动的系统也成立.从上面推导可以看出，角动量在外场的对称轴上投影总是守恒的，因为绕该轴转动时系统力学性质不变.当然，这时角动量计算是相对于位于该轴上的任意点（坐标原点）的.

最重要的情况是中心对称外场，即势能仅仅依赖于到空间中某个特定点（中心）的距离.显然，在这种场内运动时，系统角动量在任意过中心的轴上投影都守恒，就是说，系统相对场中心的角动量 $\boldsymbol{M}$ 守恒.

另一个例子是，在沿着 $z$ 轴的均匀场中角动量投影 $M_z$ 守恒，并且坐标原点可以任意选取.

应该指出，角动量在任意轴（我们就取为 $z$ 轴）上的投影，都可以由对拉格朗日函数的微分求得
\begin{equation}
  M_z=\sum_a\frac{\partial L}{\partial\dot{\varphi}_a},
\end{equation}
其中坐标 $\varphi$ 是绕 $z$ 轴的转角.根据前面给出的角动量守恒定律的证明过程，这个结论是显然的，也可以直接计算来验证.利用柱坐标 $r,\varphi,z$ 代入 $x_a=r_a\cos\varphi_a$，$y_a=r_a\sin\varphi_a$，有
\begin{equation}
  M_z=\sum_a m_a\left(x_a\dot{y}_a-y_a\dot{x}_a\right)
  =\sum_a m_a r_a^2\dot{\varphi}_a.
\end{equation}
另一方面，用这些坐标表示时，拉格朗日函数可以写成
\begin{equation*}
  L=\frac{1}{2}\sum_a m_a\left(\dot{r}_a^2+r_a^2\dot{\varphi}_a^2
  +\dot{z}_a^2\right)-U.
\end{equation*}
将它代入 (9.7) 即可得 (9.8).

\begin{xiti}

\noindent{\heiti 习题1}\quad 用柱坐标 $r$，$\varphi$，$z$ 表示质点角动量的笛卡儿坐标分量以及角动量的大小.

答案：
\begin{equation*}
\begin{aligned}
  &M_x=m\sin\varphi\,(r\dot{z}-z\dot{r})-mrz\dot{\varphi}\cos\varphi,\\
  &M_y=m\cos\varphi\,(z\dot{r}-r\dot{z})-mrz\dot{\varphi}\sin\varphi,\\
  &M_z=mr^2\dot{\varphi},\\
  &M^2=m^2r^2\dot{\varphi}^2\left(r^2+z^2\right)+m^2(r\dot{z}-z\dot{r})^2.
\end{aligned}
\end{equation*}

\noindent{\heiti 习题2}\quad 用球坐标 $r$，$\theta$，$\varphi$ 表示质点角动量的笛卡儿坐标分量以及角动量的大小.

答案：
\begin{gather*}
  M_x=-mr^2\left(\dot{\theta}\sin\varphi
  +\dot{\varphi}\sin\theta\cos\theta\cos\varphi\right),\\
  M_y=mr^2\left(\dot{\theta}\cos\varphi
  -\dot{\varphi}\sin\theta\cos\theta\sin\varphi\right),\\
  M_z=mr^2\dot{\varphi}\sin^2\theta,\\
  M^2=m^2r^4\left(\dot{\theta}^2+\dot{\varphi}^2\sin^2\theta\right).
\end{gather*}

\noindent{\heiti 习题3}\quad 在下列场中运动时动量 $\boldsymbol{P}$ 和角动量 $\boldsymbol{M}$ 的哪些分量守恒？

a. 无限大均匀平面场.

答案：$P_x,\ P_y,\ M_z$（$z$ 轴垂直于 $xy$ 平面）.

b. 无限长均匀圆柱场.

答案：$P_z,\ M_z$（$z$ 轴沿圆柱轴）.

c. 无限长均匀棱柱场.

答案：$P_z$（$z$ 轴沿棱柱）.

d. 两个点场.

答案：$M_z$（$z$ 轴过两点）.

e. 无限大均匀半平面场.

答案：$P_y$（$y$ 轴沿半平面边缘，坐标原点在边缘上）.

f. 均匀圆锥场.

答案：$M_z$（$z$ 轴沿圆锥轴）.

g. 均匀圆环场.

答案：$M_z$（$z$ 轴沿圆环轴）.

h. 无限长均匀圆柱形螺旋线场.

解：绕螺旋轴（$z$ 轴）旋转 $\delta\varphi$ 同时沿着该轴平移 $\dfrac{h}{2\pi}\delta\varphi$（$h$ 为螺距），拉格朗日函数不改变.所以有
\begin{equation*}
  \delta L=\frac{\partial L}{\partial z}\delta z
  +\frac{\partial L}{\partial\varphi}\delta\varphi
  =\left(\dot{P}_z\frac{h}{2\pi}+\dot{M}_z\right)\delta\varphi=0,
\end{equation*}
由此可得
\begin{equation*}
  P_z\frac{h}{2\pi}+M_z=\mathrm{const}.
\end{equation*}

\end{xiti}

\section{力学相似性}

拉格朗日函数乘以任意常数不会改变运动方程（参见§2）.在一些重要情况下，利用这一点，无需实际求解运动方程就可以得到有关运动性质的一些有用的结论.

这些情况包括那些势能是坐标的齐次函数的情况，即势能函数满足条件
\begin{equation}
  U(\alpha\boldsymbol{r}_1,\alpha\boldsymbol{r}_2,\cdots,\alpha\boldsymbol{r}_n)
  =\alpha^k U(\boldsymbol{r}_1,\boldsymbol{r}_2,\cdots,\boldsymbol{r}_n),
\end{equation}
其中 $\alpha$ 是任意常数，$k$ 是函数的齐次次数.

我们引入变换，使坐标都变为 $\alpha$ 倍，时间变为 $\beta$ 倍：
\begin{equation*}
  \boldsymbol{r}_a\to\alpha\boldsymbol{r}_a,\qquad t\to\beta t.
\end{equation*}
这时所有速度 $\boldsymbol{v}_a=\mathrm{d}\boldsymbol{r}_a/\mathrm{d}t$ 变为 $\alpha/\beta$ 倍，动能变为 $\alpha^2/\beta^2$ 倍，势能变为 $\alpha^k$ 倍.如果 $\alpha$ 和 $\beta$ 满足条件
\begin{equation*}
  \frac{\alpha^2}{\beta^2}=\alpha^k,
\end{equation*}
即
\begin{equation*}
  \beta=\alpha^{1-k/2},
\end{equation*}
则变换的结果是拉格朗日函数乘以常数 $\alpha^k$，运动方程保持不变.

所有质点的坐标改变相同的倍数，意味着变换前后的运动轨迹几何上相似，仅仅是尺寸不同.于是我们可以得出结论，如果系统的势能是（笛卡儿）坐标的 $k$ 次齐次函数，则由运动方程可以得到一系列几何上相似的不同轨迹，并且（不同轨迹上的相应点的）运动时间之比满足关系式
\begin{equation}
  \frac{t'}{t}=\left(\frac{l'}{l}\right)^{1-k/2},
\end{equation}
其中 $l'/l$ 是两个轨迹线度之比.除了时间以外，在相应时刻不同运动轨迹上相应点的任何力学量之比是 $l'/l$ 的幂，例如对于速度、能量和角动量有
\begin{equation}
  \frac{v'}{v}=\left(\frac{l'}{l}\right)^{k/2},\qquad
  \frac{E'}{E}=\left(\frac{l'}{l}\right)^{k},\qquad
  \frac{M'}{M}=\left(\frac{l'}{l}\right)^{1+k/2}.
\end{equation}

下面就前面所讲的举几个例子.

我们在后面章节中将会讲到，在微振动情况下势能是坐标的二次函数（$k=2$）.由 (10.2) 可知这种振动的周期与振幅无关.

在均匀力场中势能是坐标的线性函数（参见 (5.8)），即 $k=1$.由 (10.2) 得
\begin{equation*}
  \frac{t'}{t}=\sqrt{\frac{l'}{l}}.
\end{equation*}
例如，由此可知，对于重力场中的自由落体，下落时间的平方之比等于初始高度之比.

对于两个质点之间的牛顿引力或者两个电荷之间的库仑力，势能都是与两点间距离成反比，即势能是 $k=-1$ 的齐次函数.这时
\begin{equation*}
  \frac{t'}{t}=\left(\frac{l'}{l}\right)^{3/2},
\end{equation*}
例如，由此可得出结论：轨道运动周期的平方与轨道尺寸的立方成正比（这个结论称为开普勒第三定律）.

如果力学系统在有限空间中运动，势能是坐标的齐次函数，则动能和势能的时间平均值之间存在非常简单的关系，这个关系式称为位力定理.

因为动能是速度的二次函数，根据欧拉齐次函数定理有
\begin{equation*}
  \sum_a\frac{\partial T}{\partial\boldsymbol{v}_a}\cdot\boldsymbol{v}_a=2T,
\end{equation*}
或者利用 $\partial T/\partial\boldsymbol{v}_a=\boldsymbol{p}_a$ 写成
\begin{equation}
  2T=\sum_a\boldsymbol{p}_a\cdot\boldsymbol{v}_a
  =\frac{\mathrm{d}}{\mathrm{d}t}\sum_a\boldsymbol{p}_a\cdot\boldsymbol{r}_a
  -\sum_a\dot{\boldsymbol{p}}_a\cdot\boldsymbol{r}_a.
\end{equation}
我们将上面这个等式对时间平均.函数 $f(t)$ 对时间平均定义为
\begin{equation*}
  \bar{f}=\lim_{\tau\to\infty}\frac{1}{\tau}\int_0^{\tau}f(t)\,\mathrm{d}t.
\end{equation*}
容易看出，如果函数 $f(t)$ 是某个有界函数 $F(t)$ 对时间的全导数，则 $f(t)$ 对时间平均等于零.事实上，
\begin{equation*}
  \bar{f}=\lim_{\tau\to\infty}\frac{1}{\tau}\int_0^{\tau}\frac{\mathrm{d}F}{\mathrm{d}t}\,\mathrm{d}t
  =\lim_{\tau\to\infty}\frac{F(\tau)-F(0)}{\tau}=0.
\end{equation*}
假设系统在有限空间中以有限速度运动，则 $\sum_a\boldsymbol{p}_a\cdot\boldsymbol{r}_a$ 是有界的，等式 (10.4) 右端第一项对时间平均等于零.根据牛顿方程 (5.3)，将等式 (10.4) 右端第二项中 $\dot{\boldsymbol{p}}_a$ 替换为 $-\partial U/\partial\boldsymbol{r}_a$，可得\footnote{等式 (10.5) 右端表达式有时也称为系统的位力.}
\begin{equation}
  2\bar{T}=\overline{\sum_a\boldsymbol{r}_a\cdot\frac{\partial U}{\partial\boldsymbol{r}_a}}.
\end{equation}
如果势能是所有径矢 $\boldsymbol{r}_a$ 的 $k$ 次齐次函数，则根据欧拉定理，等式 (10.5) 变为所要求的关系
\begin{equation}
  2\bar{T}=k\bar{U}.
\end{equation}
由于 $\bar{T}+\bar{U}=\bar{E}=E$，可以将等式 (10.6) 等价地写成
\begin{equation}
  \bar{U}=\frac{2}{k+2}E,\qquad
  \bar{T}=\frac{k}{k+2}E,
\end{equation}
这两个公式将 $\bar{U}$ 和 $\bar{T}$ 用系统的总能量表示出来.

对于 $k=2$，即微振动的特殊情况有
\begin{equation*}
  \bar{T}=\bar{U},
\end{equation*}
即动能和势能对时间平均相等.对于牛顿引力（$k=-1$）有
\begin{equation*}
  2\bar{T}=-\bar{U}.
\end{equation*}
这时 $E=-\bar{T}$ 表明，只有在总能量为负值的情况下，在牛顿引力作用下，运动才是有界的（参见§15）.

\begin{xiti}

\noindent{\heiti 习题1}\quad 质量不同、势能相同的质点沿着相同轨道运动，它们的运动时间满足什么关系？

答案：
\begin{equation*}
  \frac{t'}{t}=\sqrt{\frac{m'}{m}}.
\end{equation*}

\noindent{\heiti 习题2}\quad 质点有相同的质量但势能相差一个常数因子，试求沿着相同轨道运动的时间之比？

答案：
\begin{equation*}
  \frac{t'}{t}=\sqrt{\frac{U}{U'}}.
\end{equation*}

\end{xiti}
